\begin{frame}{Associativity and order}
  \begin{columns}[T,onlytextwidth]
    \begin{column}{0.46\textwidth}
      \begin{exampleblock}{Associativity}
        \[
        (AB)C=A(BC)
        \]
        Grouping may change.
      \end{exampleblock}
    \end{column}
    \begin{column}{0.46\textwidth}
      \begin{alertblock}{Order}
        \[
        AB\neq BA\quad\text{in general}
        \]
        Order remains fixed.
      \end{alertblock}
    \end{column}
  \end{columns}
\end{frame}

\begin{frame}{Checkpoint: multiplication}
  \begin{columns}[T,onlytextwidth]
    \begin{column}{0.62\textwidth}
      \Large
      \begin{enumerate}
        \item check inner dimensions;
        \item predict the output shape;
        \item compute row by column;
        \item keep the order fixed;
        \item interpret the product as composition.
      \end{enumerate}
    \end{column}
    \begin{column}{0.34\textwidth}
      \takeaway{Each entry is local; the transformation is global.}
    \end{column}
  \end{columns}
\end{frame}

% =============================================================================
